\documentclass[11pt]{article}

\usepackage[margin=25mm]{geometry}

\usepackage{amsmath,amssymb,fontspec,unicode-math}

\IfFontExistsTF{DejaVu Serif}{\setmainfont{DejaVu Serif}}{\setmainfont{Latin Modern Roman}}

\IfFontExistsTF{Latin Modern Math}{\setmathfont{Latin Modern Math}}{\setmathfont{DejaVu Math TeX Gyre}}

\usepackage{xurl,hyperref}

\setlength{\parindent}{0pt}\setlength{\parskip}{6pt}

\begin{document}

\title{Mixtures of Subspaces for Bandwidth Efficient Context Parallel Training}

\author{ELVIA preliminary AI analysis}\date{21 October 2026}\maketitle

Preliminary assessment: Major.\par

Independent scientific review remains pending. Author adoption has not been requested.

Low-dimensional examples invalidate the factorized rate, local landscape correspondence and Jacobian near-isometry. Expected stochastic decay can be repaired; projection removal needs samplewise control.

\section*{Review scope}

Full supplied technical report read alongside the corresponding original-paper claim passages and applicable author-release source. This pass independently checks the stated reasoning and bounds; historical executions were not rerun.

\section*{Problem 1: Original-loss smoothness does not give the composite rate}

Source: Section 3.1; Appendix A, Lemma 1

Assessment: supported; impact: major

\subsection*{Source claim (paraphrase)}

The lemma uses original-loss smoothness for \(\Phi(B,U)=F(BUU^T)\) and simultaneous product-manifold updates.

\subsection*{Analysis}

Use \(F(W)=W_{12}^2/2\), \(U(a)=(\cos a,\sin a)^T\), \(B_0=\operatorname{diag}(2,0)\), \(a_0=0.05\), and step \(1/2\). The original loss is globally \(1\)-smooth. For \(q=b\cos a\sin a+c\sin^2a\), in \(|b-2|<0.1,|c|<0.1,|a|<0.1\), \(\partial_aq>1.9\cos0.2-0.1\sin0.2>1\), yielding PL constant \(\mu=1\). Exact simultaneous updates give \(a_1=0.05-\sin(0.2)/2\), \(B_{1,11}=1.99750832\), \(B_{1,12}=-0.00012469\). Both endpoints remain there, yet \(\Phi_1=0.00483999>\Phi_0/2=0.00249168\).

\subsection*{Finding and scope}

The claimed rate fails despite the stated smoothness, local PL and exact manifold step; the actual step still decreases loss.

\subsection*{Proposed correction}

Use a composite pullback descent constant on an invariant neighborhood. Its reciprocal step condition plus PL gives a geometric value bound; the composite constant need not equal the original one.

The described correction remains a proposal; author adoption or a corresponding erratum is not established.

\section*{Problem 2: Global coverage does not preserve local minima at folds}

Source: Section 3.2; Appendix Theorem 1

Assessment: supported; impact: major

\subsection*{Source claim (paraphrase)}

Global surjectivity or density of rotations is claimed to preserve local minima.

\subsection*{Analysis}

Let \(C=\operatorname{diag}(2,1)\), \(F(W)=\|W-C\|_F^2/2\), \(\bar U=e_1\), and \(a(s)=\pi/4+s^2+s^3\). On \(\theta\ne0\), put \(s=\log|\theta|\), \(R(\theta)=\operatorname{Rot}(a(s))\operatorname{diag}(1,\operatorname{sgn}\theta)\). This smooth map covers \(O(2)\). Minimizing over \(B\) gives \(f(a)=1/2+(3/2)\sin^2a\). Near \(\theta=1\), \(a\ge\pi/4\), so \((C,1)\) is a joint local minimum of value \(5/4\). A smaller ambient angle lowers loss, so its image is not an ambient minimum.

\subsection*{Finding and scope}

A globally covering map can fold locally; density need not even provide an exact preimage.

\subsection*{Proposed correction}

Require local openness, such as a submersion, for the reverse local-minimum implication and differential surjectivity for stationarity transfer.

The described correction remains a proposal; author adoption or a corresponding erratum is not established.

\section*{Problem 3: Rotation Jacobian stabilizers defeat near-isometry}

Source: Appendix A.1, Jacobian expansion and Corollary 1

Assessment: supported; impact: major

\subsection*{Source claim (paraphrase)}

The full rotation Jacobian is claimed near-isometric at zero, preserving Hessian conditioning.

\subsection*{Analysis}

For \(d=3,r=1,\bar U=e_1\) and Frobenius-orthonormal skew generators, \[J(0)=[e_2/\sqrt2,e_3/\sqrt2,0],\qquad J(0)^TJ(0)=\operatorname{diag}(1/2,1/2,0).\] The last generator fixes the frame. At quadratic optimum \(B=C=e_1e_1^T,U=e_1\), \(H_U=\operatorname{diag}(4,1,1)\) but \(J^TH_UJ=\operatorname{diag}(1/2,1/2,0)\), violating the claimed positive lower comparison. The skew-exponential integral derivative also gives \(\|J(\theta)\|_2\le1\), excluding the asserted forward singular-value blow-up.

\subsection*{Finding and scope}

Near-isometry fails at the clipped center. Empirical optimizer instability is not inferred.

\subsection*{Proposed correction}

Remove stabilizer coordinates and use a full-rank chart. For a positive definite Hessian, \(\kappa(J^THJ)\le\kappa(H)\kappa(J)^2\) is a valid conditional comparison.

The described correction remains a proposal; author adoption or a corresponding erratum is not established.

\section*{Problem 4: Reduced orbits can contain an optimum and create traps}

Source: Section 3.3, reduced rotation family

Assessment: supported; impact: major

\subsection*{Source claim (paraphrase)}

The no-spurious-minimum conclusion is retained if an optimum is reachable or close to the reduced orbit.

\subsection*{Analysis}

For \(d=4,r=2,\bar U=[e_1,e_3]\), \(C=\operatorname{diag}(1,0,1,0)\), rotate the two planes at speeds \(1\) and \(2\). Minimizing \(\|BUU^T-C\|_F^2/2\) over \(B\) gives \(f(\theta)=[\sin^2\theta+\sin^2(2\theta)]/2\). The orbit has zero loss at \(0\), but at \(\pi/2\), \(f=1/2,f'=0,f''=3>0\), giving a joint restricted local minimum. Rotating only the first column toward its original axis lowers ambient loss.

\subsection*{Finding and scope}

Optimum containment does not prevent restricted-orbit traps. This witness lies outside practical clipping radius \(0.5\), so it does not establish a trapped clipped run.

\subsection*{Proposed correction}

Restrict the landscape statement to the orbit or require all relevant ambient descent directions to be represented.

The described correction remains a proposal; author adoption or a corresponding erratum is not established.

\section*{Problem 5: Stochastic decay is an expected bound, not pathwise pruning}

Source: Main projection-removal box; Appendix Theorem 3

Assessment: supported; impact: minor

\subsection*{Source claim (paraphrase)}

The proof reverses the norm/expectation comparison, while the main discussion suggests a pathwise limit.

\subsection*{Analysis}

Let \(P\) be either coordinate projector equally likely, \(Q=\operatorname{diag}(1,0),\lambda=1,\eta=1/2\), and bounded smooth \(F_x(Z)=3-h(Z_{11})\), where \(h(t)=t\) on \([-2,2]\), \(|h'|\le1,|h|\le5/2\). From zero, \(W_t=w_tQ\) obeys \(w_{t+1}=w_t/2+\xi_t/2\), with independent Bernoulli halves. Thus \(Ew_t=(1-2^{-t})/2\), but arbitrarily long all-one/all-zero blocks give almost surely \(\limsup w_t=1,\liminf w_t=0\).

\subsection*{Finding and scope}

The stronger pathwise reading fails; the formal expected bound survives a corrected proof. Derivatives through trainable projectors remain a separate condition.

\subsection*{Proposed correction}

Take triangle inequality before expectation to obtain \(y_{t+1}\le(1-\eta\lambda)y_t+\eta Lp_Q\), sum the geometric series, and state an expected limsup.

The described correction remains a proposal; author adoption or a corresponding erratum is not established.

\section*{Problem 6: Average overlap alone does not justify projection removal}

Source: Section 3.5, projection removal

Assessment: supported; impact: minor

\subsection*{Source claim (paraphrase)}

Weight decay is used to infer redundancy of sample-dependent projections.

\subsection*{Analysis}

Let coordinate projectors \(P_1,P_2\) be equally likely, \(F_i(Z)=1-e^{-\|Z-P_i\|_F^2}\), and \(\lambda=e^{-1/4}\). Gradients are globally bounded by \(\sqrt{2/e}<1\). The regularized objective has Hessian lower bound \((e^{-1/4}-2e^{-3/2})I\succ0\) and zero gradient at \(W_*=I/2\), hence this is its unique optimum. For each sample, removal changes loss from \(1-e^{-1/4}\) to \(1-e^{-1/2}\), an increase \(0.17227012\). No globally idle direction exists.

\subsection*{Finding and scope}

The overlap theorem lacks samplewise residual control even at an optimum. The reported removal/fine-tuning experiment remains separate empirical evidence.

\subsection*{Proposed correction}

Expected residual \(E\|XW(I-P(x))\|_F\le\varepsilon\) and downstream Lipschitz constant \(K\) suffice for expected loss change at most \(K\varepsilon\); the overlap bound does not establish this premise.

The described correction remains a proposal; author adoption or a corresponding erratum is not established.

\section*{Limits}

Preliminary AI source/argument assessment; independent human scientific adjudication remains pending.\par

No author-provided or retained verification program was executed. Published empirical results were not independently reproduced.\par

\section*{Source identity}

Authors’ arXiv:2606.16384v1, current official version checked 2026-10-21

\url{https://arxiv.org/pdf/2606.16384v1}

Reviewed PDF SHA-256: \nolinkurl{3e7e4da1d9d9c98201e8a17bff13a68c98f36e5ace9a88023601357addc31659}

The SHA-256 identifies the exact current-author-version PDF inspected. The original acquisition timestamp for this retained PDF is unknown; file modification times were not used. Current version identity was independently checked against official arXiv metadata.

\end{document}
